\section{Automating the Exploit}
\label{sec:automation}

The Bash harvester iterated over competition teams at
\code{fd66:666:<team>::2}, TCP/1337, excluding our own box and
infrastructure targets. It checked whether registration responded before
attempting the chain. For each reachable target it:

\begin{enumerate}
  \item registered a short-lived account whose ship name contained the
        file-listing command, then triggered its \code{typ} launcher;
  \item registered another carrier for the grep command and triggered it;
  \item fetched the output through the \code{Path} traversal;
  \item normalized and deduplicated flags locally, then submitted them
        through MOTH.
\end{enumerate}

Short one-character usernames saved bytes in the constrained launcher.
Each registration also produced a session cookie needed to add the
component. The form-encoded launcher used \code{\%20} for spaces.

MOTH received batches tagged with the service and source team. Its bearer
token came from \code{MOTH\_API\_TOKEN} in the environment, rather than a
credential embedded in the public script. Team jobs ran concurrently with
a worker limit so that one slow instance did not hold up the entire sweep.

Teams patched while the competition ran. The harvester had to distinguish
unreachable services, broken registrations, rejected primitives, and
instances where the full chain still worked. A failed request alone did
not establish which of those states applied.
